%% toc: P15 161 161.3,161.12
\section{Observation fields smaller than the sample space}
\label{sec:annex-p15-w02c-p15-017-observation-fields}

\begin{annexitem}{annex:p15-w02c-p15-017-observation-field-proper}{thm:p15-w02c-p15-017-observation-field-proper}{A rectangle field that contains every observation}
\begin{proof}
\emph{Part 1.} The set \(\Omega=\Omega'\times\{0,1\}\) belongs to
\(\mathcal G\), with \(A'=\Omega'\). For \(A'\in\mathcal F'\),
\[
 \Omega\setminus(A'\times\{0,1\})=(\Omega'\setminus A')\times\{0,1\},
\]
and for \(A'_1,A'_2,\ldots\in\mathcal F'\),
\[
 \bigcup_k\bigl(A'_k\times\{0,1\}\bigr)
 =\Bigl(\bigcup_kA'_k\Bigr)\times\{0,1\}.
\]
Both right sides lie in \(\mathcal G\), because \(\mathcal F'\) is a
sigma-algebra. Hence \(\mathcal G\) is a sigma-algebra. Each of its
members is a measurable rectangle, so \(\mathcal G\subseteq\mathcal F\).

For a measurable value set \(D\in\mathcal S_n\),
\[
 X_n^{-1}(D)=\{(\omega',y):X'_n(\omega')\in D\}
 =X_n'^{-1}(D)\times\{0,1\}\in\mathcal G,
\]
since \(X'_n\) is measurable. Thus \(\sigma(X_n)\subseteq\mathcal G\) for
every \(n\). By \cref{def:p15-tail-sigma-algebra},
\(\mathcal H_\infty\) is the sigma-algebra generated by the union of the
fields \(\sigma(X_n)\); since \(\mathcal G\) is a sigma-algebra containing
that union, \(\mathcal H_\infty\subseteq\mathcal G\). Finally
\(\mathcal T\subseteq\mathcal F'_1=\sigma(X_1,X_2,\ldots)=\mathcal H_\infty\).

\emph{Part 2.} The set \(B=\Omega'\times\{1\}\) is a measurable
rectangle, so \(B\in\mathcal F\). Suppose \(B=A'\times\{0,1\}\) for some
\(A'\in\mathcal F'\). The probability space \(\Omega'\) is nonempty,
since \(P'(\Omega')=1\); choose \(\omega'\in\Omega'\). Then
\((\omega',1)\in B\), so \(\omega'\in A'\), and therefore
\((\omega',0)\in A'\times\{0,1\}=B\). This contradicts the definition of
\(B\), whose points all have second coordinate one. Hence
\(B\notin\mathcal G\). Since \(\mathcal H_\infty\subseteq\mathcal G\), also
\(B\notin\mathcal H_\infty\), while \(B\in\mathcal F\); so
\(\mathcal H_\infty\neq\mathcal F\). No property of the \(X'_n\) was used.

\emph{Part 3.} Let \(A\in\mathcal H_\infty\). By part~1,
\(A=A'\times\{0,1\}\) for some \(A'\in\mathcal F'\). Then
\[
 A\setminus B=A'\times\{0\},\qquad
 B\setminus A=(\Omega'\setminus A')\times\{1\},
\]
and these two rectangles are disjoint. The product measure gives
\[
 P(A\triangle B)
 =\tfrac12P'(A')+\tfrac12\bigl(1-P'(A')\bigr)=\tfrac12 .
\]
Also \(A\cap B=A'\times\{1\}\), so
\(P(A\cap B)=\tfrac12P'(A')\), while \(P(A)=P'(A')\) and
\(P(B)=\tfrac12\). Hence \(P(A\cap B)=P(A)P(B)\). Since
\(P(A\triangle B)=1/2\neq0\) for every \(A\in\mathcal H_\infty\), no event
of \(\mathcal H_\infty\) coincides with \(B\) outside a null set.

\emph{Part 4.} For each \(n\), \(\mathcal H_n=\sigma(X_1,\ldots,X_n)\) is
generated by a subfamily of the generators of \(\mathcal H_\infty\), so
\(\mathcal H_n\subseteq\mathcal H_\infty\). Part~3 then gives
\(P(B\triangle C)=1/2\) for every \(C\in\mathcal H_n\). The conclusion of
\cref{thm:p15-finite-coordinate-approximation} for the event \(B\) and
any \(\varepsilon\leq1/2\) would require some \(C\in\mathcal H_n\) with
\(P(B\triangle C)<1/2\), and no such \(C\) exists. So that conclusion is
false for \(B\in\mathcal F\setminus\mathcal H_\infty\).

For the example's smallest case, let \(\Omega'=\{\omega'\}\). Then
\(\mathcal F'=\{\emptyset,\Omega'\}\), so
\(\mathcal G=\{\emptyset,\Omega\}\). Every sigma-algebra contains
\(\emptyset\) and \(\Omega\), so part~1 gives
\(\mathcal T=\mathcal H_\infty=\{\emptyset,\Omega\}\). The rectangle field
on the two-point set \(\Omega\) is its power set, with four members.
\end{proof}
\end{annexitem}

\begin{annexitem}{annex:p15-w02c-p15-017-canonical-generated}{thm:p15-w02c-p15-017-canonical-generated}{Coordinate generators}
\begin{proof}
For each \(i\geq1\) and \(C\in\mathcal S\),
\(X_i^{-1}(C)=\{\omega:\omega_i\in C\}\). Hence the union
\(\bigcup_i\sigma(X_i)\) is exactly the family of generating sets of the
coordinate product field \(\mathcal F\), because
\(\sigma(X_i)=\{X_i^{-1}(C):C\in\mathcal S\}\) (the preimages under one
map of the sets of a sigma-algebra form a sigma-algebra, and every
member of \(\sigma(X_i)\) is such a preimage). By
\cref{def:p15-tail-sigma-algebra},
\(\mathcal H_\infty=\sigma\bigl(\bigcup_i\sigma(X_i)\bigr)\), which is the
sigma-algebra generated by the same family as \(\mathcal F\). So
\(\mathcal H_\infty=\mathcal F\). The argument used only the form of
the generators, so it applies unchanged when the \(i\)th coordinate
ranges over its own measurable space \((S_i,\mathcal S_i)\).

The preimage family is a sigma-algebra because
\(X_i^{-1}(S)=\Omega\),
\(\Omega\setminus X_i^{-1}(C)=X_i^{-1}(S\setminus C)\), and
\(\bigcup_kX_i^{-1}(C_k)=X_i^{-1}(\bigcup_kC_k)\). It contains every
generator of \(\sigma(X_i)\) and is contained in any sigma-algebra that
makes \(X_i\) measurable, so it equals \(\sigma(X_i)\).
\end{proof}
\end{annexitem}
